%! TeX program = lualatex
\documentclass[10pt,a4paper,table]{article}
\usepackage[english]{babel}
\usepackage{amsmath}
\usepackage{fontspec}
\setmainfont{Times New Roman}
\newfontfamily\quotefont[Ligatures=TeX]{Linux Libertine O}
\newfontfamily\vestb{TRY Vesterbro Poster}
\newfontfamily\overp{Overpass}
\usepackage{xcolor}
\definecolor{frbgreenl}{RGB}{162,202,137}
\definecolor{frbgreen}{RGB}{28,85,66}
\usepackage{graphicx}
\graphicspath{{../img/}}
\usepackage{tikz}
\usetikzlibrary{calc,shadows}
\usepackage[pdfborder={0 0 0}]{hyperref}
\usepackage{enumitem}
\usepackage{titlesec}
\setlength{\parskip}{0.5\baselineskip}
\setlength{\parindent}{0cm}
\setlist[enumerate]{label=\alph*),left=20pt}
\newcommand{\qm}[1]{``#1''}
\renewcommand\thesubsection{Exercise \arabic{subsection}}
\titleformat{\subsection}[leftmargin]{\bfseries}{}{0pt}{\hspace{-5em}\thesubsection}
\usepackage{fancyhdr}
\fancypagestyle{plain}{%
  \setlength{\headheight}{60pt}%
  \setlength{\headsep}{0pt}
  \fancyhf{}
  \renewcommand{\footrulewidth}{0pt}%
  \renewcommand\headrule{}%
  \fancyhead[R]{\overp\color{frbgreenl}\normalsize English\\\huge\color{frbgreen}\vestb Be Right Back}
  \cfoot{\thepage}
}
\pagestyle{fancy}
\fancyhf{}
\renewcommand\headrule{}
\cfoot{\thepage}

\begin{document}
\thispagestyle{plain}
\begin{tikzpicture}[remember picture,overlay]
  \node[anchor=north west,yshift=-48pt,xshift=120pt,opacity=0.2]
    at (current page.north west)
    {\includegraphics[scale=0.25]{frbvuc_logo.png}};
\end{tikzpicture}
\begin{tikzpicture}[remember picture,overlay]
  \definecolor{bgline}{RGB}{205,214,208}
  \coordinate (C) at ($(current page.north)+(35mm,-5mm)$);
  \draw[bgline,line width=0.6pt] (C) circle (65mm);
  \draw[bgline,line width=0.4pt]
    ($(current page.north east)+(0,0mm)$) ++(-28mm,0)
    -- ($(current page.south east)+(0,10mm)$) ++(-28mm,0);
\end{tikzpicture}

\begin{enumerate}
  \item In groups, select and analyse a scene from the film that you particularly liked, found important or were intrigued by. In your analysis, you must include a focus on filmic devices such as close-ups, perspective, sound effects or lighting. Moreover, connect the scene to the episode as a whole and give reasons why you picked this scene.

  \item Split up the groups from the previous task. Individually, walk about in the classroom and listen to analyses of scenes made by students from the other groups.

  \item Would you consider the episode soft or hard sci-fi? Give reasons for your answer.

  \item Identify the two most important turning points in the story and give reasons for your choices. Compare your choices to your classmates'.

  \item Characterise Martha and discuss if she is a flat or complex character.

  \item Explain the dramatic irony when Martha's sister tells her that she is glad Martha \qm{is moving on}.

  \item Using your notes from \textit{Focus for your watching}, characterise Ash and the replica of Ash and compare the two.

  \item Compare Martha's relationship with the real Ash to her relationship with the replica. Find arguments that the new relationship is better as well as worse and find scenes to support your answers.

  \item The Ash replica keeps reminding Martha that he is not real and she is fully aware of this all the way through. What is it then that initially makes her accept the replica as a substitute and later reject him?

  \item When Martha and the Ash replica are at the cliff, \qm{Ash} twice says that her sentences are difficult to process. Find the exact wording of what Martha says and discuss what it is that makes these sentences difficult for artificial intelligence.

  \item Discuss if Martha ought to have let her sister in on her secret. Might the relationship with the replica have turned out differently then?

  \item Revisit Isaac Asimov's Three Laws of Robotics and discuss if the Ash replica follows all the rules. Are there scenes where the laws are tested and, if so, what is the role of these scenes?

  \item Bearing the YouTube video from the pre-reading in mind, discuss how the film is an extrapolation of contemporary issues and technology.

  \item Discuss the ending of the episode:
    \begin{itemize}
      \item Is it an open ending or is there a resolution?
      \item Is it a happy ending? Why/why not?
      \item Does Martha achieve closure? If so, how? If not, why not?
      \item How does the ending with the attic form a connection to the beginning and what is the effect of this?
      \item In conclusion, what are the main themes and what do we learn about them?
    \end{itemize}
\end{enumerate}

\end{document}
